\documentclass[preprint,12pt]{elsarticle}

\usepackage{booktabs}
\usepackage{array}
\usepackage{url}
\usepackage{amsmath}
\usepackage{enumitem}

\emergencystretch=2em

\journal{Cyber Security and Applications}

\begin{document}

\begin{frontmatter}

\title{Recovery-State Semantic Observability and Decision Identifiability Across Public Space-Cyber and Real-Flight Data Products}

\author[aff1]{Aman Kumar Singh\corref{cor1}}
\ead{asingh65430@ucumberlands.edu}
\cortext[cor1]{Corresponding author. ORCID: 0009-0008-9752-3743.}
\affiliation[aff1]{organization={Independent Researcher},
            city={The Woodlands},
            state={Texas},
            country={United States}}

\begin{abstract}
Public satellite cybersecurity datasets support anomaly detection and attack analysis, but recovery decisions require semantically valid evidence of trust, freshness, epoch, contradictions, evidence completeness, security signals, and authorization. We integrate two independently governed investigations using the same frozen eight-variable recovery-state interface, while preserving distinct populations and estimands. Study 9 evaluated CuCD-ID v3, AegisSat 2025, and UNSW-IoTSAT 2026 under a preregistered mapping rubric and four frozen downstream policies. All three datasets had 0/8 operational direct or permitted derived coverage. Their original row populations of 25,000, 137,965, and 404,798 collapsed losslessly to one decision-equivalent group each; none yielded a uniquely identifiable action. Exact guaranteed minimum additional sidecar-state sizes were six, seven, seven, and eight variables across the four policies. Study 10 independently assessed published data-product schemas for ESA Anomaly Dataset v2 Missions 1--3 and OPS-SAT-AD v2. Its 32 descriptive source-variable classifications comprised 14 ambiguous and 18 absent mappings, with no direct or predeclared-derived mappings. Because the published products did not establish a native recovery decision point and evidence-validity window, the 16 stratum-policy operational action assessments were correctly not evaluable, rather than assigned an action rate. In a within-family sensitivity comparison, all 224 ESA channels matched between releases, while 54 of 896 comparable metadata cells differed and v2 added a \texttt{Categorical} column. Synthetic oracle parity covered 1,024 complete policy-input comparisons without mismatches. The integration was decided after both studies' results and is not a single prospective experiment. Together, these findings delimit the semantic information released for recovery-oriented research; they do not assess spacecraft recovery efficacy or operational policy safety.
\end{abstract}

\begin{keyword}
satellite cybersecurity \sep semantic observability \sep cyber resilience \sep spacecraft telemetry \sep recovery state \sep decision identifiability
\end{keyword}

\end{frontmatter}

\section{Introduction}

Public satellite cybersecurity datasets have become increasingly useful for developing and evaluating intrusion-detection methods, attack classifiers, telemetry analytics, and cyber-physical testbeds. CuCD-ID provides labelled CubeSat command and telemetry data generated through a NOS3/cFS environment \cite{fayyaz2026cucd}; AegisSat couples an Earth-based CubeSat with an environment emulator and attack manager \cite{idan2025aegissat}; and UNSW-IoTSAT provides a large hybrid testbed-derived smart-satellite dataset spanning cyber, physical, and radio-frequency features \cite{abdelhameed2026unsw}. These resources broaden access to repeatable space-cybersecurity experimentation.

Detection-oriented data, however, are not automatically response-ready data. A recovery controller may need to know whether evidence is cryptographically valid, whether a source belongs to a trusted set, whether evidence is fresh for the current decision epoch, whether accepted evidence is contradictory, whether a minimum evidence set is complete, whether an operational security signal is present, and whether an authorization path is available. Intrusion-response research distinguishes the response decision from detection alone \cite{bashendy2023irs}, while cyber-resiliency engineering emphasizes the ability to recover from and adapt to adverse conditions as part of a broader systems-security objective \cite{ross2021cyberresilient}. The semantic interface between security telemetry and recovery logic therefore matters independently of detector accuracy.

Cross-dataset cybersecurity research commonly focuses on predictive generalization, feature harmonization, or label normalization. Distribution shift can degrade performance even when in-distribution results are strong \cite{koh2021wilds}, and recent IoT work has shown both the difficulty of cross-network transfer \cite{cerasuolo2026transfer} and the value of explicit rule-based label harmonization \cite{rodriguez2026harmonization}. The present study addresses a different question. It does not test whether a classifier transfers between datasets. Instead, it asks whether independently generated public space-cyber datasets expose the semantics required by a fixed downstream recovery interface and, when they do not, whether the downstream recovery action remains identifiable.

The study was preregistered as S9-RTSI-001. Three public datasets were frozen before row-level endpoint inspection, a field-level semantic mapping rubric was frozen before canonical execution, and four downstream recovery policies were evaluated separately. Missing state remained unknown rather than being inferred from attack labels, lexical similarity, or analyst expectation. The resulting negative finding is itself informative: none of the three datasets established any of the eight required recovery-state variables directly or through a permitted deterministic derivation, and no frozen policy produced a uniquely identifiable recovery action from native dataset state alone.

The contribution is therefore not a new detector, a new dataset, or a claim of operational spacecraft recovery. It is a reproducible semantic-interface analysis that (i) distinguishes field presence from recovery meaning, (ii) propagates unresolved semantics into downstream action identifiability, and (iii) computes the minimal additional sidecar state needed to remove decision ambiguity without imputing unavailable mission or trust semantics.


\paragraph{Integration scope.}
This manuscript reports two separately frozen studies with different evaluation populations and units. Study 9 (\texttt{S9-RTSI-001}) provides field-level semantic mapping, admissible-completion policy analysis, and exact minimal-sidecar results for three public space-cyber datasets. Study 10 (\texttt{S10-RSEI-001}) asks whether the same recovery-state interface is visible in four descriptive source strata from two released real-flight anomaly-detection products, ESA Anomaly Dataset v2 \cite{esaad2025v2} and OPS-SAT-AD v2 \cite{opssat2024v2}. Study 10 is an external applicability investigation, not a post-hoc addition to Study 9's frozen population. The decision to integrate their manuscripts was made after the individual outcomes were known; it is a publication decision, not a retrospectively preregistered combined analysis.

\subsection{Research questions}

The study follows four preregistered research questions:

\begin{description}[leftmargin=1.15cm,labelwidth=0.95cm,style=nextline]
    \item[RQ1] Across preregistered public space-cyber datasets, which recovery decision variables are directly observable, deterministically derivable by a predeclared rule, ambiguous, or absent without using attack labels as operational truth?
    \item[RQ2] Conditioned only on qualifying dataset-native operationally visible state, when is the downstream recovery action uniquely identifiable versus ambiguous across admissible completions of missing decision state?
    \item[RQ3] What minimal additional sidecar or interface state is sufficient to make the downstream recovery action uniquely identifiable without imputing unavailable mission, trust, freshness, epoch, contradiction, or authorization semantics?
    \item[RQ4] Which missing recovery-state semantics are common across independently generated public space-cyber datasets and which are dataset-specific?
\end{description}

The Results section answers these questions in order through semantic coverage, policy-stratified action identifiability, guaranteed minimal sidecar state, and the cross-dataset common unresolved set.

\section{Related Work and Research Gap}

\subsection{Public space-cybersecurity datasets and testbeds}

CuCD-ID provides labelled CubeSat cybersecurity data generated in NOS3/cFS with COSMOS reproduction tooling \cite{fayyaz2026cucd}. Its raw population contains 25,000 records. AegisSat was developed as an open satellite cybersecurity testbed combining a physical CubeSat, environmental emulation, controlled attacks, and labelled telemetry from repeated experiments \cite{idan2025aegissat}. UNSW-IoTSAT extends the smart-satellite dataset landscape with a hybrid cyber-physical testbed and more than 404,000 labelled records containing satellite, sensor, communications, and RF-related features \cite{abdelhameed2026unsw}.

These datasets were not created to instantiate the recovery-state interface evaluated here. Their inclusion therefore does not imply an expectation that they should contain every recovery variable. The purpose of the comparison is to measure semantic sufficiency under a fixed interface, not to grade dataset quality.

\subsection{Cross-dataset transfer and harmonization}

Cross-dataset generalization is a known challenge in machine-learning systems. WILDS formalized evaluation under naturally occurring distribution shifts and showed that in-distribution performance does not guarantee out-of-distribution robustness \cite{koh2021wilds}. In cybersecurity, cross-network evaluation similarly shows that intrusion-detection models can lose effectiveness when transferred to heterogeneous environments \cite{cerasuolo2026transfer}. Rule-based harmonization has also been proposed to transform heterogeneous IoT security datasets into more interoperable feature and label representations \cite{rodriguez2026harmonization}. Cybersecurity ontology research likewise uses semantic interoperability to unify heterogeneous defensive knowledge sources \cite{akbar2023ontology}.

Study 9 is deliberately narrower. It does not harmonize attack taxonomies, propose semantic interoperability as a general concept, or train a transferable model. It tests whether the native operational fields of each selected dataset semantically establish a preregistered recovery-state contract. Statistical correlation with a missing recovery variable is not treated as semantic observability.

\subsection{Intrusion response and cyber resilience}

Intrusion-response systems add a decision layer after detection, selecting countermeasures or actions based on the state available to the response mechanism \cite{bashendy2023irs}. NIST cyber-resiliency guidance similarly frames recovery and adaptation as systems-engineering concerns that depend on trustworthy system state and risk-informed decisions \cite{ross2021cyberresilient}. Space-specific guidance also treats response and recovery as explicit engineering functions: NIST IR 8441 applies the Cybersecurity Framework to hybrid satellite networks and emphasizes interfaces among independently operated segments while retaining Detect, Respond, and Recover functions \cite{mccarthy2023hsn}. Sandia's Adaptive Intrusion Response for Space Systems (AIRSS) work likewise addresses space-system cyber response by classifying cyberattacks and identifying mitigation responses in real time \cite{sahakian2021airss}.

These foundations motivate the central distinction in this paper: response selection itself is not the contribution. AIRSS and related response work ask which mitigation should be selected once a threat state is available, whereas Study 9 asks whether independently produced public space-cyber datasets expose enough semantically valid recovery state for a fixed downstream recovery action to be identifiable without imputation.

The contribution therefore does not rest on claiming that semantic interoperability, space-system cyber response, or recovery are individually new. The research gap addressed here is their narrower conjunction: field-level semantic mapping across independently sourced public space-cyber datasets, exhaustive downstream action-identifiability analysis with missing state preserved as unknown, and minimal-sidecar derivation without using attack labels as operational truth.

\section{Study Design}

\subsection{Two-study design and nonexchangeable units}
Study 9 retains its original frozen RQ1--RQ4, three datasets, finite source-row denominators, four-policy action sets, and exact minimal-sidecar analysis. Study 10 retains its separately frozen RQ-E1--RQ-E5, four descriptive source strata (three ESA v2 missions and one OPS-SAT-AD v2 stratum), schema-level mapping endpoints, conditional action-evaluability rule, synthetic controls, and optional matched ESA v1/v2 release sensitivity. Both use the pinned eight-variable Study 9 semantic contract and the four Study 2 policy identifiers, without altering either. Scientific outcomes were generated under separate protocols. Study 9 source rows cannot be pooled with Study 10 source-variable slots, and the ESA missions are not treated as independent samples of spacecraft effectiveness. The post-results manuscript integration is disclosed here explicitly.

\subsection{Frozen recovery-state interface}

The downstream interface contains eight binary recovery-state variables:
\begin{enumerate}[label=(\arabic*)]
    \item \texttt{signature\_valid},
    \item \texttt{source\_trusted},
    \item \texttt{fresh},
    \item \texttt{epoch\_valid},
    \item \texttt{contradictory},
    \item \texttt{minimum\_evidence\_complete},
    \item \texttt{security\_signal}, and
    \item \texttt{authorization\_available}.
\end{enumerate}

The frozen Study 2 selector is used only as a bounded downstream decision interface. It is not retrained, optimized, or revalidated by this study. The selector identity is fixed by SHA-256 in the Study 9 governance records.

\subsection{Dataset population}

The preregistered primary population consists of CuCD-ID v3, AegisSat 2025, and UNSW-IoTSAT 2026. Each dataset is analyzed as a separate finite source population; no pooled inferential population is constructed. Canonical source-row counts are shown in Table~\ref{tab:datasets}.

\begin{table}[htbp]
\centering
\caption{Frozen Study 9 dataset populations.}
\label{tab:datasets}
\begin{tabular}{lrr}
\toprule
Dataset & Source rows & Native-state groups \\
\midrule
CuCD-ID v3 & 25,000 & 1 \\
AegisSat 2025 & 137,965 & 1 \\
UNSW-IoTSAT 2026 & 404,798 & 1 \\
\bottomrule
\end{tabular}
\end{table}

\subsection{Semantic mapping rubric}

For dataset $d$ and required recovery variable $v$, the frozen mapping function assigns
\[
M_d(v) \in \{\mathrm{DIRECT},\mathrm{DERIVABLE},\mathrm{AMBIGUOUS},\mathrm{ABSENT}\}.
\]
\texttt{DIRECT} requires a native operational field whose documented semantics establish the recovery variable. \texttt{DERIVABLE} requires a frozen deterministic rule whose inputs are themselves available under the allowed interface. \texttt{AMBIGUOUS} means one or more fields are semantically related but do not establish the required proposition. \texttt{ABSENT} means no eligible native field represents the proposition.

Attack or scenario labels are excluded as operational recovery evidence. They may describe the source population, but they cannot be converted into trust, authorization, freshness, evidence completeness, contradiction state, cryptographic validity, or an operational security signal for the primary endpoint.

No deterministic derivation rules were permitted in the frozen Study 9 adjudication. Consequently, direct-or-derivable coverage can exceed direct coverage only if a prospectively frozen derivation exists; none did.

\subsection{Protocol deviation for UNSW \texttt{Position\_Anomaly}}

The initial semantic adjudication treated \texttt{Position\_Anomaly} as a direct binary security signal. Full-population validation found 72 of 404,798 canonical rows outside the documented 0/1 domain. Before accepting canonical results, a protocol deviation reclassified the field from direct to ambiguous for the entire frozen UNSW population. The analysis did not drop the violating rows, salvage only binary-looking rows, coerce non-zero values, realign neighboring columns, or substitute attack labels. This conservative correction was recorded before the accepted canonical run.

\subsection{Decision identifiability}

Let $U_d$ be the unresolved recovery-state variables for dataset $d$. For each native record or losslessly collapsed native-state group $g$, all binary assignments to $U_d$ are treated as admissible completions. For policy $p$, the reachable action set is
\[
A_{d,p}(g)=\{f_p(c):c\in C_d(g)\},
\]
where $f_p$ is the frozen selector under policy $p$ and $C_d(g)$ is the set of admissible completions. A recovery action is uniquely identifiable only when
\[
|A_{d,p}(g)|=1.
\]
The row-level unique-action fraction is computed with exact multiplicity weighting after lossless state collapse.

Four policies were frozen and evaluated separately:
\begin{itemize}[nosep,leftmargin=*]
    \item \texttt{S2\_B0\_FAIL\_CLOSED},
    \item \texttt{S2\_B1\_FAIL\_OPERATIONAL},
    \item \texttt{S2\_B2\_RISK\_THRESHOLD}, and
    \item \texttt{S2\_S1\_EVIDENCE\_AWARE}.
\end{itemize}
Policy pooling or averaging is prohibited.

\subsection{Guaranteed minimal sidecar}

For unresolved set $U$, a sidecar subset $S\subseteq U$ is sufficient only if revealing every possible assignment to $S$ guarantees a singleton reachable action set across every assignment to the unrevealed variables $U\setminus S$. The endpoint retains every sufficient subset of minimum cardinality. Actual unavailable values are never substituted merely to reduce ambiguity.

\subsection{Finite-population reporting and audit}

Primary reporting uses deterministic finite-population counts, fractions, set operations, and exact combinatorial summaries. No sampling model was introduced, so no $p$-values or confidence intervals are reported. The canonical implementation uses lossless decision-equivalent state collapse with exact multiplicity preservation. A separate independent audit reconstructs raw-row projection and policy-independent coverage without calling the canonical projector or mapping implementation. Accepted outputs are bound by SHA-256 in the frozen result record.


\subsection{Study 10: real-flight published-schema observability (RQ-E1--RQ-E5)}
Study 10 evaluated the v2 ESA Anomaly Dataset for Missions 1--3 \cite{esaad2025v2} and OPS-SAT-AD v2 \cite{opssat2024v2}. The frozen primary units were four descriptive source strata times eight recovery-state variables, yielding 32 classification slots, \emph{not} 32 independent experiments or 32 independent spacecraft observations. Study 10 applied the same fixed classes \texttt{DIRECT}, \texttt{DERIVABLE\_BY\_PREDECLARED\_RULE}, \texttt{AMBIGUOUS}, and \texttt{ABSENT}, with an empty derivation registry and a required visibility role for each classification. A field was operationally direct only when a documented native field supplied the decision-relevant value at the required decision point. Source identifiers did not establish trusted membership; timestamps without a decision reference and validity interval did not establish freshness; retrospective anomaly labels did not count as live security signals; and telecommand presence did not establish an available authorized contact path.

Study 10's bounded source inventory checked nine version-pinned local files, reviewed 12 ESA published metadata tables and both OPS-SAT CSV headers, and confirmed the ESA v2 outer-directory organization of 224 channel and 821 telecommand archives. It did not deserialize the nested pickled telemetry tables or inspect individual samples. The ESA outer-layout JSON check is supported by the author's terminal-reported digest rather than independently received report bytes; the G4-05 boundary and G4-06 schema JSON reports were independently rehashed. The primary outcome is therefore limited to the documented \emph{published data-product representation}, not hidden spacecraft capabilities or unseen flight processing.

Under the frozen conditional action rule, Study 10 could enumerate admissible policy action sets only for a defensible native recovery decision time and evidence-validity window. If that prerequisite was absent, the specified result was \texttt{NOT\_EVALUABLE\_FOR\_OPERATIONAL\_ACTION}, not a zero rate or an assumed false recovery-state vector. Separately, synthetic G3 controls tested 18 prespecified cases and all 256 complete Boolean inputs under each of four policies against an independent oracle; these are methodology tests and do not contribute to primary empirical coverage.

The optional Study 10 paired sensitivity compared ESA v1 \cite{esaad2024v1} and v2 \cite{esaad2025v2} \texttt{channels.csv} metadata by exact (mission, channel identifier, field) matching. It counted shared-field equality/difference and release-specific columns without reading nested channel telemetry. ESA v1 was not a fifth primary stratum. Study 10's source-rights status remained conditional under the locally authorized R1 evidence standard, and no third-party dataset contents were redistributed.

\section{Results}

\subsection{Recovery-state coverage (RQ1)}

Table~\ref{tab:coverage} summarizes the mapping classes. No dataset contains a variable classified as \texttt{DIRECT} or \texttt{DERIVABLE}. CuCD-ID has four ambiguous and four absent variables. AegisSat and UNSW-IoTSAT each have five ambiguous and three absent variables.

\begin{table}[htbp]
\centering
\caption{Recovery-state semantic coverage under the frozen rubric.}
\label{tab:coverage}
\begin{tabular}{lrrrr}
\toprule
Dataset & Direct & Derivable & Ambiguous & Absent \\
\midrule
CuCD-ID v3 & 0 & 0 & 4 & 4 \\
AegisSat 2025 & 0 & 0 & 5 & 3 \\
UNSW-IoTSAT 2026 & 0 & 0 & 5 & 3 \\
\bottomrule
\end{tabular}
\end{table}

Accordingly, operational direct coverage is 0/8 for every dataset, and operational direct-or-derivable coverage is also 0/8 for every dataset. The cross-dataset common unresolved set is the complete eight-variable recovery interface.

The mapping is not based on field-name absence alone. For example, source identifiers in all three datasets are related to provenance but do not establish membership in a trusted source set. Time and sequence fields exist in CuCD-ID and AegisSat but do not provide the decision reference and validity semantics required for freshness or epoch validity. UNSW-IoTSAT includes CRC, quality, satellite-identity, timestamp, and position-anomaly fields, yet those fields do not establish cryptographic signature validity, trusted-source membership, recovery evidence completeness, or a population-wide boolean security signal under the frozen definitions.

\subsection{Lossless native-state collapse}

After semantic projection, all eight required variables remain unresolved for every row in each dataset. Therefore each source population forms a single decision-equivalent native-state group. This collapse is computationally lossless: multiplicities of 25,000, 137,965, and 404,798 are retained exactly and used in all row-level denominators. It is not sampling, deduplication, or row exclusion.

\subsection{Policy-stratified action identifiability (RQ2)}

For all three datasets and all four policies, the unique-action fraction is 0. Table~\ref{tab:policies} reports the reachable action sets and guaranteed sidecar cardinalities. Because every dataset projects to the same unresolved eight-variable interface, the policy-level action sets are identical across datasets; only the source-row multiplicities differ.

\begin{table}[htbp]
\centering
\small
\caption{Policy-level reachable actions and guaranteed minimal sidecar size. Human-readable action labels are used here; exact frozen action identifiers are retained in the canonical result artifacts.}
\label{tab:policies}
\begin{tabular}{>{\raggedright\arraybackslash}p{0.18\linewidth}>{\raggedright\arraybackslash}p{0.56\linewidth}c}
\toprule
Policy & Reachable actions from native state & Sidecar size \\
\midrule
B0 fail-closed &
Restrict and request authorization; proceed to recovery gate & 6 \\
B1 fail-operational &
Preserve limited operation; restrict and request authorization & 7 \\
B2 risk-threshold &
Hold and require evidence; preserve limited operation; restrict and request authorization & 7 \\
S1 evidence-aware &
Hold and require evidence; restrict and request authorization; proceed to recovery gate & 8 \\
\bottomrule
\end{tabular}
\end{table}

For the fail-closed policy, the guaranteed minimal sidecar contains six variables: signature validity, source trust, freshness, epoch validity, contradiction state, and minimum evidence completeness. The fail-operational and risk-threshold policies additionally require the security signal, giving seven variables. The evidence-aware policy requires all eight variables, including authorization availability.

\subsection{Guaranteed minimal sidecar state (RQ3)}

The exact guaranteed minimum is policy-dependent rather than dataset-dependent because all three datasets project to the same unresolved eight-variable recovery interface. The minimum cardinalities are 6 for B0 fail-closed, 7 for B1 fail-operational, 7 for B2 risk-threshold, and 8 for S1 evidence-aware. These values answer RQ3 under the frozen selector and do not imply that the same number of variables is universally required by other recovery architectures.

\subsection{Cross-dataset common missing semantics (RQ4)}

The intersection of unresolved recovery state across the three datasets equals the entire required interface:
\begin{itemize}[nosep,leftmargin=*]
    \item \texttt{signature\_valid},
    \item \texttt{source\_trusted},
    \item \texttt{fresh},
    \item \texttt{epoch\_valid},
    \item \texttt{contradictory},
    \item \texttt{minimum\_evidence\_complete},
    \item \texttt{security\_signal}, and
    \item \texttt{authorization\_available}.
\end{itemize}
No dataset-specific missing variable remains after taking this intersection because all eight are unresolved in all three datasets.


\subsection{Study 10: published-data-product observability (RQ-E1 and RQ-E2)}
The separate Study 10 primary analysis classified eight recovery variables within each of four descriptive release strata. Table~\ref{tab:study10coverage} records the complete frozen-schema mapping totals. All 32 slots were assigned as 14 ambiguous and 18 absent. None was a qualifying operational direct field or a permitted predeclared derivation. The resulting \emph{published-schema-level} direct and direct-or-derivable coverage was 0/8 within each stratum. These four descriptive outcomes are not an additional row-level population of Study 9, and a seven-stratum pooled coverage percentage is not calculated.

\begin{table}[htbp]
\centering
\caption{Study 10 classification counts by published real-flight data-product stratum. These are schema-level descriptions, not independent flight trials.}
\label{tab:study10coverage}
\begin{tabular}{lrrrr}
\toprule
Source stratum & Direct & Derivable & Ambiguous & Absent \\
\midrule
ESA v2 Mission 1 & 0 & 0 & 4 & 4 \\
ESA v2 Mission 2 & 0 & 0 & 4 & 4 \\
ESA v2 Mission 3 & 0 & 0 & 3 & 5 \\
OPS-SAT-AD v2 & 0 & 0 & 3 & 5 \\
\bottomrule
\end{tabular}
\end{table}

The ESA products contain channel and anomaly metadata, with telecommand records present in the documented Mission 1 and 2 layouts; the OPS-SAT release contains timestamped telemetry-related fields and derived anomaly-detection features. Their released representations do not contain a documented, complete recovery evidence state conforming to the pinned interface. An \texttt{ABSENT} mapping means absent from the documented published representation under the adjudication contract; it does not assert that the mission never implemented the underlying onboard function. An \texttt{AMBIGUOUS} mapping means related metadata do not uniquely establish the required decision proposition. Neither may be coerced to a known false Boolean.

\subsection{Study 10: conditional action evaluability (RQ-E4)}
Across four Study 10 strata and four frozen policy identifiers, the predeclared result is \texttt{NOT\_EVALUABLE\_FOR\_OPERATIONAL\_ACTION} for all 16 reporting combinations. Published-schema evidence did not establish the required native recovery decision point and evidence-validity window. No operationally grounded action set, singleton-action fraction, row-level denominator, or recoverability rate was calculated. This outcome must be distinguished from Study 9's exact mathematical exhaustive-completion analysis: the Study 9 zero unique-action fractions are \emph{not} operational flight-recovery failure rates, and they are not substituted as Study 10 outcomes.

\subsection{Study 10: within-family ESA release sensitivity (RQ-E3)}
The optional matched release analysis identified the same 224 channel identifiers in ESA v1 and v2: 76 Mission 1, 100 Mission 2, and 48 Mission 3, with zero v1-only or v2-only channels. Of 896 comparable cells in the four shared fields \texttt{Subsystem}, \texttt{Physical Unit}, \texttt{Group}, and \texttt{Target}, 842 were equal and 54 differed. All 54 differences were in \texttt{Physical Unit}: zero in Mission 1, ten in Mission 2, and 44 in Mission 3. The new v2 \texttt{Categorical} metadata field was present for all 224 channels and is counted as a release-specific field addition, \emph{not} 224 changed existing cells. Table~\ref{tab:esareleases} retains the paired within-mission reporting.

\begin{table}[htbp]
\centering
\caption{Matched ESA channel-metadata sensitivity. Comparisons concern annotations, not nested telemetry payloads or recovery semantics.}
\label{tab:esareleases}
\begin{tabular}{lrrrr}
\toprule
Mission & Matched & Equal & Changed & New in v2 \\
\midrule
ESA Mission 1 & 76 & 304 & 0 & 76 \\
ESA Mission 2 & 100 & 390 & 10 & 100 \\
ESA Mission 3 & 48 & 148 & 44 & 48 \\
\midrule
Total & 224 & 842 & 54 & 224 \\
\bottomrule
\end{tabular}
\end{table}

Changed physical-unit metadata does not, without the original unit values and an authoritative explanation, establish a unit conversion, a semantic change, an error correction, or a different series payload. No nested channel stream was opened, and no matched-release result is counted as independent spacecraft replication.

\subsection{Study 10: synthetic technical parity (RQ-E5)}
The frozen synthetic-control implementation passed all 18 cases and 22 repository-native unit tests in the recorded GitHub Actions execution. An independently coded four-policy oracle matched the pinned Study 2 selector for 1,024 comparisons (256 complete Boolean inputs per policy), with zero mismatches. This validates the defined synthetic interface and its implementation parity, not real cryptographic verification, operational authorization, or spacecraft recovery effectiveness.

\section{Discussion}


\subsection{Cross-study synthesis and the boundary of external applicability}
Both frozen studies find zero operational direct recovery-state coverage under a shared eight-variable semantic interface, but they answer different questions. Study 9 establishes the consequences of unresolved native state \emph{for a fully defined mathematical downstream selector}: exact admissible-action sets and minimum guaranteed additional sidecar state under four frozen policies. Study 10 establishes a different boundary: released real-flight anomaly-detection products do not document a complete decision-time recovery evidence state suitable for grounding the selector operationally. The latter finding cannot be expressed as Study 9's action-identifiability rate without manufacturing a recovery decision point, and the former cannot be presented as empirical spacecraft recovery performance.

The incremental role of Study 10 is therefore an external \emph{representation-level} applicability assessment and release-version sensitivity, not a new demonstration of policy effectiveness. The shared zero coverage must not be framed as seven independent replications or as a pooled estimate of space-system resilience. The ESA source family contains three curated descriptive mission strata, and independence of those missions from the OPS-SAT lineage has not been independently established. Identical high-level coverage outcomes do not establish identical native semantics, data-generating mechanisms, or security capabilities.

\subsection{Post-results integration and research integrity}
Study 9 and Study 10 retained distinct frozen protocols, scopes, preexisting source exposures, analysis units, outcomes, evidence pins, and execution authorizations. The decision to integrate them into this manuscript was made after their respective results were available. The combined manuscript is therefore \emph{not} a jointly preregistered study, and no Study 9 population expansion or retrospective alteration of the Study 10 design is claimed. The publication integration was undertaken to synthesize related evidence transparently while avoiding two overlapping accounts of the same recovery-state interface.

\subsection{Limits of the evidence added by Study 10}
The Study 10 dataset inventory covered metadata tables, OPS-SAT CSV headers, and ESA outer archive directory layouts plus published format documentation, not the contents of every nested flight-series object. The published-data-product mapping cannot rule out undocumented capabilities or support operational performance conclusions. The author-local original-source integrity check and ESA layout receipt have different verification provenance; the layout JSON digest was reported in terminal output rather than independently checked against uploaded bytes. ESA/OPS-SAT source-use rights and notices are conditionally qualified; third-party redistribution rights are not assumed. The ESA paired comparison summarizes metadata differences without exposing channel metadata values, so its 54 changed values cannot be interpreted as verified changes in physical units or telemetry meanings.

\subsection{Meaning of the negative result}

The principal result is a semantic non-identifiability finding, not a dataset-quality judgment. The analyzed datasets contain substantial telemetry, attack labels, timing information, identifiers, counters, quality indicators, and experimental context. What they do not provide, under the frozen Study 9 contract, is enough explicit recovery-state meaning to select a unique downstream action without additional state.

This distinction is important because presence of a related field is weaker than establishment of a recovery proposition. An application identifier is not the same as a trusted-source decision. A timestamp is not equivalent to freshness unless a decision-time reference and validity rule are defined. A CRC result is not a cryptographic signature verification result. A command field does not establish that an authorization channel is available. Preserving those distinctions prevents a response system from silently inheriting assumptions that were never represented in the source data.

\subsection{Implications for dataset and interface design}

If a cybersecurity dataset is intended only for detection or attack classification, omission of recovery semantics may be entirely appropriate. If the same data are expected to support intrusion response, resilience analysis, or recovery-policy evaluation, the findings suggest that an explicit sidecar interface should accompany the telemetry.

The sidecar need not reproduce the entire spacecraft state. It should expose the propositions required by the decision interface with clear operational semantics and provenance. In the frozen selector studied here, the minimum guaranteed requirement ranges from six to eight binary variables depending on policy. That result provides a concrete interface-design target rather than a generic recommendation to collect ``more data.''

\subsection{Semantic interoperability versus predictive transfer}

The study also separates two forms of interoperability. Predictive transfer asks whether a learned model trained in one domain performs in another. Semantic interoperability asks whether two data sources represent propositions with equivalent decision meaning. The two can fail independently. A model may transfer statistically even when the target dataset lacks an explicit recovery-state interface, and two datasets may expose compatible semantics even when their feature distributions differ.

This distinction complements existing cross-dataset work rather than replacing it. Distribution-shift benchmarks \cite{koh2021wilds}, cross-network intrusion-detection studies \cite{cerasuolo2026transfer}, and rule-based harmonization frameworks \cite{rodriguez2026harmonization} address model and label portability. Study 9 isolates the downstream decision semantics that remain after detection output is available.

\subsection{Protocol conservatism}

The UNSW \texttt{Position\_Anomaly} deviation illustrates why semantic claims must be population-wide and fail closed. The field was strongly related to the intended security-signal concept and documented as binary, but 72 released rows violated that domain. Selectively dropping or repairing those rows would have produced a more favorable direct-coverage result at the cost of changing the frozen population. Reclassifying the entire field as ambiguous preserved the released artifact and the preregistered semantics.

Likewise, attack labels were not promoted into operational state even though doing so would have reduced ambiguity. Offline ground truth answers a different question from what a controller knows at decision time.

\subsection{Limitations}

Several limitations bound interpretation. First, only three frozen public space-cyber datasets are evaluated; the result is not representative of every satellite cybersecurity dataset or operational mission. Second, semantic sufficiency is defined relative to one frozen eight-variable recovery interface and four frozen selector policies. Different response architectures may require different propositions. Third, the study does not estimate intrusion-detection accuracy, classifier transferability, attack prevalence, or causal mission effects. Fourth, no flight hardware, on-orbit recovery, mission operator, or authorization authority validates the downstream actions. Fifth, the finite-population analysis intentionally makes no statistical generalization beyond the frozen artifacts.

The result also does not establish that six to eight sidecar variables are universally necessary for satellite recovery. Those cardinalities are exact only for the frozen selector, policies, and unresolved interface examined here.

\section{Reproducibility and Data Availability}

The accepted Study 9 science is frozen in the public repository at
\url{https://github.com/Zartharas/mission-aware-satellite-cyber-recovery}.
The canonical run is identified as
\texttt{S9-RTSI-001\_CANONICAL\_RUN\_20260914} under \texttt{study9/results/}.
The immutable results-freeze record retains the complete SHA-256 identities of the accepted outputs. For concise manuscript traceability, the accepted science merge and governance-closeout commits begin with \texttt{1dfdef4ec9a0} and \texttt{ced258fc3ba9}, respectively; their complete identifiers remain in the repository governance records.

The repository contains protocol, schema, mapping, code, tests, deterministic outputs, independent-audit outputs, and artifact identity records. Third-party raw dataset bytes are not redistributed in the repository; the verification records bind the analyzed public artifacts by source identity and cryptographic hash.


The Study 10 protocol, frozen design receipt, and G3 implementation evidence are maintained under \texttt{study10/}. The completed G4 Tier A/C and ESA paired-release outputs are presently retained as hash-identified author research packages outside the canonical GitHub manuscript tree; this draft does not claim that those G4 outputs have been deposited in a stable public repository. The source-specific releases are ESA Anomaly Dataset v2 \cite{esaad2025v2}, OPS-SAT-AD v2 \cite{opssat2024v2}, and ESA Anomaly Dataset v1 solely for paired sensitivity \cite{esaad2024v1}. The Study 10 result package was prepared as a manuscript-integration input after the outcome calculations, with provenance and non-reproduction boundaries recorded. Both study protocols and frozen source artifacts remain unchanged; no raw ESA or OPS-SAT telemetry is redistributed. Dataset credit, applicable licenses, changes, and non-endorsement/notice requirements must be retained in final submission materials. The exact accepted study-specific hashes and source attestations are enumerated in the governance records rather than inferred from this combined text.

\section{Conclusion}

Two independently governed investigations examined what publicly released space-cybersecurity and real-flight telemetry data products make observable for a fixed eight-variable recovery-decision interface. Study 9 found zero operationally direct or predeclared-derived coverage in CuCD-ID v3, AegisSat 2025, and UNSW-IoTSAT 2026. Its frozen downstream selector produced no uniquely identifiable action from native state alone under any of four policies, and exact guaranteed additional sidecar requirements ranged from six to eight variables. Study 10 independently mapped four published ESA and OPS-SAT data-product strata and found 14 ambiguous and 18 absent classifications, again with zero qualifying operational direct or derived fields. Unlike Study 9, its operational action-identifiability endpoint was not evaluable without a native decision point and evidence-validity interval. A bounded release comparison matched all 224 ESA channels but detected 54 shared-field metadata differences and a new 224-entry v2 categorical column; these observations do not establish changes in flight telemetry or recovery semantics.

These complementary outcomes delineate a boundary between anomaly-detection data and recovery-decision evidence. Where recovery-oriented evaluation is a goal, source publications and operational interfaces would need to expose explicit trust, temporal validity, evidence completeness, contradiction, security-guard, and authorization semantics. The results do not establish actual spacecraft recovery performance, security-control absence, or universal policy superiority. The integration of the two studies is a disclosed post-results publication decision.

\section*{CRediT authorship contribution statement}

\textbf{Aman Kumar Singh:} Conceptualization; Methodology; Software; Validation; Formal analysis; Investigation; Resources; Data curation; Writing -- original draft; Writing -- review and editing; Visualization; Project administration.

\section*{Funding}

This research was conducted independently and received no specific grant from funding agencies in the public, commercial, or not-for-profit sectors.

\section*{Declaration of competing interest}

The author declares no competing financial or non-financial interests.

\section*{Declaration of generative AI and AI-assisted technologies in the manuscript preparation process}

OpenAI ChatGPT was used to support manuscript drafting, language refinement, structural review, and consistency checking against the separately frozen Study 9 and Study 10 research records and integrated manuscript evidence. The author reviewed and edited the resulting text, verified scientific claims against the frozen repository evidence, and assumes full responsibility for the manuscript.

\bibliographystyle{elsarticle-num}
\bibliography{references}

\end{document}
